\section{Conclusion}
\label{sec:conclusion}

We presented AgenticRL, a trajectory-aware controller for a frozen multimodal
RAG stack. A budgeted refinement loop couples generation-visible retained
evidence $M^+$ with controller-only failed retrieval traces $M^-$, while
generator-derived utility computed after executed retrieval trains only the
controller.

On EndoBench, cross-round state is what makes budgeted iterative retrieval pay
off: the same controller and frozen stack reach 39.56\% without it, below the
40.12\% single-round agent, but 42.84\% with it ($+3.28$ points,
$p<10^{-12}$), improving the closed-book generator by 2.25 points and correcting
10.66\% of single-round errors. The protocol-matched ablation attributes
$+1.30$ of the matched gain to the state itself, the gain is confined to the
questions on which the state becomes active, and the state simultaneously
reduces retrieval calls, parse failures, and empty-evidence fallbacks under
identical execution rules. The resulting system, a 4B controller over a frozen
8B generator, is 1.15 points above the 41.69\% GPT-4o accuracy reported by
EndoBench and also exceeds its 34B and 80B medical-specialist models, with
gains concentrated on spatial and procedural reasoning.

Retained evidence pays off only when a later round can reuse it, and its
isolated matched contribution is positive but not significant, so the strongest
evidence supports the complete fixed-budget, evidence-preserving strategy rather
than a protocol-independent memory effect. The reward requires gold answers
during training and is most natural for multiple-choice supervision. The
controller also inherits the behavior of its frozen retrieval environment, and
iterative retrieval adds latency. Future work will study open-ended rewards,
transfer across indexes and generators, cost-aware stopping, and prospective
evaluation; AgenticRL remains a retrieval research framework rather than a
clinical decision system.
